\section{Conclusion}

This study investigated SLMs ability to generate code and analyzed how many test cases the model passes by verifying the execution of the generated code, to better understand the main types of errors in these models. We conducted experiments to verify whether SLMs are capable of generating correct code, what the most frequent errors are based on the degree of error, the ability of the models to identify these errors in the code, how specialization can improve this task, and how the models can be compared with static analyses and what the impact of these analyses is on the evaluation and generation of code by the models.

The code-generation results show that SLMs can generate valid code, producing operational code in more than 90\% of generations and reaching almost 80\% fully functional generation. However, the main error of the models lies in the inability to understand all the nuances of functional requests, which leads to deviations from the requested tasks, resulting in a \emph{Functional Error}. The least frequent type of hallucination is \emph{Syntax Error}, which highlights the model's difficulty in following the language's grammar. In one case, there were many errors; however, these were not due to code errors but to the model using natural language instead of the requested language. These results corroborate previous work pointing to this behavior; however, this experiment shows that models still struggle with it, although they reduce it across all hallucination types. This analysis also used feedback generated by static tools, which reduced the number of codes with \emph{Functional Error}; however, it affected other classes, including those that should not be affected by the feedback, and some performed worse on the compilable code response.

The results on the models' ability to detect code hallucinations, using a process called LLM-as-a-judge, show a divergence between code generation and evaluation, since models that perform well at code generation are not necessarily good at code review. The main difficulty with these models is identifying Runtime Errors, meaning that while they can generate code with correct execution flows, they may not detect certain problems. However, compared with static tools, which can detect runtime errors, SLM analysis is, in most cases, superior.

A strength of some models was their ability to detect functional errors in the analyzed code, demonstrating an interesting skill as reviewers of requested functionality, something static tools cannot measure, making these local models a promising use case. However, although the models can identify these flaws, the correction is not always effective; although the models that managed to self-detect managed to improve the accuracy rate, the conversion was not total.

Using static tools in the cycle, along with a self-evaluation mechanism, allowed some models to correct invalid writing errors and, in all cases, runtime errors. However, in some models, iterativity was less stable, with cases of regression. Regression is also evident when the task is limited to evaluating code for runtime errors. Because Pyright performs worse than the models, the models began to treat its output as absolute truth and, even with the warning not to, ended up relying on static analysis and missing some cases. The effective use of the added context was for \emph{Syntax Error}, as most models managed to have complete detection in most scenarios with this feedback, which reinforces that, although the models do not have the capabilities to detect syntax errors during reading, integration with other tools can be useful, especially within an iterative cycle, so that the model provides code free of this type of error.

Finally, the specialized hallucination-detection model, while showing superior results compared to smaller traditional models, still performs significantly worse. Some considerations relate to the training target: due to the unbalanced dataset, the trained models showed low capacity to detect lower-representation models and demonstrated a bias toward selecting the majority class during training.

\subsection{Future Work}

The most important direction for future work is to evaluate how these models behave in environments that replicate real-world software development. In real-world development, logic isn't tied to a single function; it spans the entire work environment, including files, external APIs, dependencies, local requirements, and specific execution requirements. A natural extension is the use of benchmarks that investigate at the repository level. However, some problems don't map as effectively as in the three classes we've presented. Adding the complexity of repository-level work introduces other considerations, such as context control and engineering, and models' ability to use tools in an agentic loop. These are variables that should be considered in a later study.

Another advantage to leverage in real-world development environments is verifying other runtime problems. In other analyzed works, many errors consisted of hallucinations of repository symbols, often due to a lack of context. By adding real-world development environments and using an agentic, interactive cycle, it's possible to observe more deeply what types of execution errors models commit, even when the context includes the symbol to be used.

Other works can perform less complex operations, which would expand the current analysis by measuring the effect of using these models without quantization on these results. Furthermore, for model training, it is interesting to expand the use of variable elements, from using other models for training to using a larger training set. However, this approach can be computationally very costly and time-consuming.